% FFI26 — Study: bảng tổng hợp tài liệu (literature matrix), tách khỏi main.tex.
% Biên dịch: cd report && latexmk -xelatex literature_matrix.tex
% Nguồn: docs/study.md (đọc bài ngày 31/10/2026; vị trí §/Eq./Table ghi trong từng ô).
% FPO, π0-FPO, DSRL, LIBERO, Ng 1999: đọc bản HTML arXiv ngày 02/11/2026, chưa có trong docs/study.md.
% Số SmolVLA lấy từ ffi26/ffi26/results/probe_report.json (viết tắt "probe").
\documentclass[9pt,a4paper]{extarticle}

\usepackage{fontspec}
\usepackage{polyglossia}
\setmainlanguage{vietnamese}
\setmainfont{Latin Modern Roman}
\setsansfont{Latin Modern Sans}
\setmonofont{Latin Modern Mono}

\usepackage[landscape,margin=1.4cm,top=1.5cm,bottom=1.5cm]{geometry}
\usepackage{amsmath,amssymb}
\usepackage{booktabs,tabularx,array,longtable}
\usepackage{enumitem}
\usepackage{xcolor}
\usepackage{caption}
\usepackage{float}
\usepackage[hidelinks]{hyperref}

\setlist{itemsep=0pt,topsep=1pt,leftmargin=1em}
\captionsetup{font=small,labelfont=bf,skip=4pt}
\setlength{\parskip}{2pt}
\renewcommand{\arraystretch}{1.25}
\newcolumntype{Y}{>{\raggedright\arraybackslash}X}
\newcolumntype{L}[1]{>{\raggedright\arraybackslash}p{#1}}
\newcommand{\code}[1]{\texttt{\detokenize{#1}}}
\newcommand{\review}[1]{\fcolorbox{orange}{orange!8}{\parbox{0.97\linewidth}{\scriptsize\textcolor{orange!70!black}{\textbf{[Cần bạn điền / kiểm tra]} #1}\\[1.6em]}}}

\title{\textbf{Study: bảng tổng hợp tài liệu}\\[2pt]\large RL cho VLA tổng quát; diffusion và flow matching trong sinh hành động}
\date{}

\begin{document}
\maketitle
\vspace{-2.5em}


% ===================================================================== 1 RL cho VLA
\section*{1. Phương pháp RL cho VLA và policy flow}
{\scriptsize
\begin{longtable}{@{}L{2.0cm} L{2.6cm} L{3.6cm} L{2.7cm} L{3.0cm} L{3.0cm} L{3.3cm} L{3.4cm}@{}}
\caption{Ma trận so sánh các bài RL cho VLA và policy flow.}\label{tab:rl}\\
\toprule
Bài & Model nền, dạng action & Mục tiêu học & Reward & Cách huấn luyện, phần được update & Giả định & Kết quả chính & Câu hỏi còn mở \\
\midrule\endfirsthead
\toprule
Bài & Model nền, dạng action & Mục tiêu học & Reward & Cách huấn luyện, phần được update & Giả định & Kết quả chính & Câu hỏi còn mở \\
\midrule\endhead
\bottomrule\endfoot
\textbf{SimpleVLA-RL} 
& OpenVLA-OFT 7B, head sửa thành token; \textbf{token}
& GRPO, không critic, bỏ KL; clip $[0{,}8;\,1{,}28]$; bỏ nhóm toàn thắng/toàn thua; nhiệt độ rollout 1,6
& nhị phân 0/1 theo kết quả, gán cho mọi token của trajectory 
& online, sim; SFT trước rồi RL; 8$\times$A800; update toàn bộ
& SFT phải có success $>0$; reset được môi trường
& LIBERO TB 91,0$\to$99,1, SFT đầy đủ; SFT 1 traj./task: 48,9$\to$96,9, Long 17,3$\to$91,7 
& học khi success $\approx0$; policy học đường tắt khác demo, đẩy vật thay vì gắp, ``pushcut''\\
\midrule
\textbf{$\pi_{\text{RL}}$}
& $\pi_0$, $\pi_{0.5}$, 3,3B; \textbf{flow}
& PPO + GAE ở mức chunk; log-likelihood qua Flow-Noise hoặc Flow-SDE
& nhị phân; ManiSkill thêm 0,1 khi gripper chạm vật
& online, sim, khung RLinf; SFT few-shot trước; chỉ update action expert $\sim$300M, VLM đóng băng
& có SFT trước; sim đủ nhanh để rollout online
& LIBERO few-shot $\pi_0$: 57,6$\to$97,6 (Flow-Noise) / 96,1 (Flow-SDE); $\pi_{0.5}$: 77,1$\to$98,3 / 97,9
& cần nhiều rollout và reset; Flow-Noise hay Flow-SDE tốt hơn tuỳ model; MetaWorld OOD không cải thiện \\
\midrule
\textbf{$\pi^*_{0.6}$ / RECAP} 
& $\pi_{0.6}$: Gemma 3 4B + action expert 860M; \textbf{flow} (+ token FAST phụ)
& advantage conditioning: $-\log\pi(a\mid o,\ell)-\alpha\log\pi(a\mid I,o,\ell)$, $I=\mathbb{1}(A>\epsilon_\ell)$; value phân phối 201 bin, VLM 670M; không cần log-prob chính xác
& $-1$ mỗi bước, 0 nếu thành công, $-C_{\text{fail}}$ nếu thất bại
& lặp offline trên robot thật: thu rollout + can thiệp của người, huấn luyện lại từ checkpoint pretrain; update toàn bộ VLA + value function
& quan sát Markov; can thiệp của người luôn tốt ($I=$ True); có người gán reward, can thiệp, reset
& robot thật: throughput tăng hơn 2 lần, tỉ lệ thất bại giảm khoảng một nửa ở task khó; success $>90\%$ trừ task giặt đồ đa dạng
& phụ thuộc người reset, gán reward, can thiệp \\
\midrule
\textbf{RIPT-VLA} 
& QueST $\sim$20M (token); OpenVLA-OFT 7B (hồi quy liên tục); \textbf{token / liên tục}
& PPO clip + advantage leave-one-out (RLOO), $K=16$ (QueST) / 8 (OFT); bỏ context toàn thắng/toàn thua
& nhị phân 0/1, không critic
& online, sim; OFT cần thêm head dự đoán scale Laplace để có log-prob; OFT: LoRA $r=32$ + action head
& SFT có success $>0$; reset được đúng context ban đầu cho cả nhóm rollout
& LIBERO, QueST: 82,7$\to$93,6, Long 68,8$\to$87,5; OFT: 96,7$\to$97,5; 1 demo: 3,5\%$\to$97,2\%
& học khi success $\approx0$; SFT đã mạnh thì RL tăng ít (OFT +0,8) \\
\midrule
\textbf{VLA-RL}
& OpenVLA-7B; \textbf{token}
& PPO + GAE, có critic, critic warmup; log-prob = tổng log-prob của 7 token action
& sparse của môi trường + reward model (RPRM) quá trình học từ nhãn giả : tách bước con theo độ mở gripper và vận tốc $\approx0$
& online, sim (LIBERO 4 suite, mỗi suite train riêng); LoRA
& nhãn giả từ gripper/vận tốc khớp ranh giới bước con
& LIBERO TB 76,5$\to$81,0, Long 53,7$\to$59,8
& RPRM dựa trên heuristic, tác giả thừa nhận khó áp dụng cho thao tác khéo léo \\
\midrule
\textbf{FPO} 
& không dùng model pretrain: policy flow train từ đầu (MuJoCo Playground, humanoid); \textbf{flow}
& PPO-clip, thay tỉ lệ likelihood bằng $\exp(\hat{\mathcal L}_{\text{CFM},\theta_{\text{old}}}-\hat{\mathcal L}_{\text{CFM},\theta})$; $\epsilon_{\text{clip}}=0{,}05$; không cần log-prob chính xác
& reward của môi trường (DeepMind Control; humanoid theo DeepMimic)
& online, sim; update policy + value function như PPO
& loss flow matching bằng $-$ELBO cộng hằng số, hằng số triệt tiêu trong tỉ lệ
& MuJoCo TB 759,3 so với PPO Gauss 667,8; humanoid chỉ theo root: 54,3\% so với 29,8\%
& chưa thử trên policy pretrain bằng BC; thiếu công cụ như ước lượng KL\\
\midrule
\textbf{$\pi_0$-FPO}
& $\pi_0$ \textbf{đóng băng}; một actor flow nhỏ sinh latent, $\pi_0$ giải mã latent thành action; \textbf{flow}
& FPO: chênh lệch loss flow matching, chuẩn hoá theo batch, $\rho=\exp(\beta z)$, rồi PPO-clip; critic ensemble + GAE
& reward thưa của môi trường, không mô tả chi tiết
& online, sim (LIBERO 4 suite, ALOHA); chỉ update actor latent + critic
& giảm loss flow matching tương ứng tăng mật độ, chưa chứng minh
& LIBERO TB 87,2, Long 65,3 so với $\pi_0$-FAST 85,5 / 60,2
& không báo seed, phương sai, siêu tham số; không sửa trọng số $\pi_0$ nên gần DSRL hơn là fine-tune VLA \\
\midrule
\textbf{DSRL}
& policy diffusion/flow bất kỳ, kể cả $\pi_0$ 3,3B; \textbf{flow / diffusion}
& học policy chọn nhiễu đầu vào $w$, coi base policy là một phần môi trường; SAC hoặc DSRL-NA (2 critic)
& reward của môi trường; robot thật dùng 0/1
& online, offline hoặc offline-to-online; sim và robot thật; base policy \textbf{đóng băng}, chỉ cần chạy forward
& base policy xử lý được nhiễu khác nhau cho hành vi khác nhau
& $\pi_0$ trên LIBERO: $\approx$20\%$\to\approx$100\% sau $\approx$10\,000 mẫu; robot thật, $\pi_0$ bật máy nướng bánh: 5/20$\to$18/20
& bị giới hạn bởi hành vi base policy làm được; chưa biết trước mức cải thiện tối đa\\
\end{longtable}}

\paragraph{Câu hỏi còn mở chung} (tổng hợp của em từ cột cuối):
\begin{enumerate}[leftmargin=1.5em]
\item \textbf{Reward thưa.} Reward nhị phân không cho gradient khi không rollout nào thành công (SimpleVLA-RL, RIPT-VLA); reward quá trình (VLA-RL) lại dựa trên heuristic.
\item \textbf{Hiệu quả mẫu và reset.} RL online cần nhiều rollout và reset về cùng cảnh ban đầu; trên robot thật phải cần người ($\pi^*_{0.6}$).
\item \textbf{Log-likelihood cho policy flow.} Chưa có cách tốt nhất cho mọi model ($\pi_{\text{RL}}$: Flow-Noise hay Flow-SDE tuỳ model); RECAP và FPO tránh log-prob; DSRL tránh cả gradient qua base policy.
\end{enumerate}

% ===================================================================== 2 Diffusion vs flow
\section*{2. Diffusion và flow matching trong sinh hành động}
{\scriptsize
\begin{longtable}{@{}L{3.0cm} L{10.8cm} L{10.8cm}@{}}
\caption{Ma trận so sánh diffusion (Diffusion Policy, DDPM) và flow matching ($\pi_0$, SmolVLA).}\label{tab:gen}\\
\toprule
Tiêu chí & Diffusion (DDPM) — Diffusion Policy, arXiv:2303.04137 & Flow matching — $\pi_0$ arXiv:2410.24164; SmolVLA arXiv:2506.01844 \\
\midrule\endfirsthead
\toprule
Tiêu chí & Diffusion (DDPM) & Flow matching \\
\midrule\endhead
\bottomrule\endfoot
\textbf{Điểm chung}
& \multicolumn{2}{L{21.8cm}}{Một mạng nhận \emph{chunk action đang bị nhiễu} cùng bước thời gian/mức nhiễu, có điều kiện là quan sát, và trả ra một đại lượng (nhiễu hoặc vận tốc) để khử nhiễu dần từ Gauss về action.} \\
\midrule
\textbf{Kiến trúc}
& Điều kiện là quan sát $O_t$. Hai backbone: CNN 1D với FiLM theo $O_t$ và bước $k$; hoặc transformer, token action nhiễu đi qua cross-attention với embedding quan sát.
& VLM xử lý ảnh, câu lệnh và state (prefix). \emph{Action expert} là transformer nhỏ xử lý token action và attend vào đặc trưng VLM; KV của prefix được cache nên mỗi bước tích phân chỉ tính lại phần action ($\pi_0$). \newline
SmolVLA (probe): VLM SmolVLM2-500M (16 lớp ngôn ngữ, vision 12 lớp); tổng 450,0M tham số, action expert 98,2M (16 lớp, rộng 0,75 lần VLM); prefix 177 token (2 camera $\times$ 64 + 48 ngôn ngữ + state); output chunk 50$\times$7; expert cross-attention vào VLM, xen self-attention mỗi 2 lớp. \\
\midrule
\textbf{Mục tiêu huấn luyện}
& Lấy action sạch $x^0$, chọn bước $k$, cộng nhiễu $\varepsilon^k$; mạng \emph{dự đoán nhiễu}:
$\mathcal{L}=\mathrm{MSE}\big(\varepsilon^k,\ \varepsilon_\theta(O_t,\,x^0+\varepsilon^k,\,k)\big)$.
& Nội suy thẳng giữa action và nhiễu; theo quy ước code SmolVLA $x_t=t\varepsilon+(1-t)A$, $t=1$ là nhiễu thuần (probe; bài $\pi_0$/SmolVLA viết $\tau$ ngược chiều: $A^\tau=\tau A+(1-\tau)\varepsilon$). Mạng \emph{dự đoán vận tốc}: $\big\|v_\theta(x_t,o)-(\varepsilon-A)\big\|^2$ (SmolVLA). $t$ lấy theo phân phối Beta, ưu tiên mức nhiễu cao ($\pi_0$: $\mathrm{Beta}((s-\tau)/s;\,1{,}5,\,1)$, $s=0{,}999$). \\
\midrule
\textbf{Inference}
& Từ nhiễu Gauss $x^K$, lặp $K$ bước $x^{k-1}=\alpha\big(x^k-\gamma\,\varepsilon_\theta(x^k,k)+\mathcal{N}(0,\sigma^2I)\big)$; $\alpha,\gamma,\sigma$ theo noise schedule. Train 100 bước; inference DDIM 10 bước, 0,1\,s trên RTX 3080.
& Lấy nhiễu một lần, tích phân ODE bằng Euler từ $t=1$ về 0: $x\leftarrow x-\Delta t\,v_\theta$. $\pi_0$ và SmolVLA dùng 10 bước ($\pi_0$, $\delta=0{,}1$; probe \code{num_steps = 10}). SmolVLA: 96,5\,ms/chunk (batch 1, RTX 5070 Ti; probe). \\
\midrule
\textbf{Đường từ nhiễu đến action}
& theo noise schedule $\alpha,\gamma,\sigma$
& đường thẳng (nội suy tuyến tính) \\
\midrule
\textbf{Tính ngẫu nhiên khi sinh}
& ở mọi bước (DDPM cộng nhiễu mỗi bước)
& chỉ ở nhiễu ban đầu; phần còn lại tất định (ODE) \\
\midrule
\textbf{Log-likelihood của action}
& mỗi bước DDPM là một chuyển Gauss, nên log-prob của cả chuỗi khử nhiễu tính được
& ODE tất định: log-likelihood chính xác phải tích phân divergence của $v_\theta$ dọc đường đi (như continuous normalizing flow), rất tốn; loss flow chỉ là cận dưới ($\pi^*_{0.6}$) \\
\midrule
\textbf{Hệ quả cho RL}
& policy gradient (PPO, GRPO) dùng trực tiếp log-prob theo từng bước khử nhiễu
& cần một trong các hướng: (1) chèn nhiễu học được vào mỗi bước Euler, mỗi bước thành Gauss ($\pi_{\text{RL}}$ Flow-Noise); (2) đổi ODE thành SDE giữ nguyên phân phối biên, thành MDP hai tầng ($\pi_{\text{RL}}$ Flow-SDE); (3) bỏ log-prob, điều kiện policy theo advantage (RECAP); (4) xấp xỉ tỉ lệ likelihood bằng chênh lệch loss flow matching (FPO). SmolVLA cùng cấu trúc với $\pi_0$ nên (1), (2) áp dụng được bằng cách thêm nhiễu vào vòng Euler và thêm critic. \\
\end{longtable}}

% ===================================================================== 3 Nền
\section*{3. Tài liệu nền}
\begin{itemize}
\item \textbf{Reward shaping} Thêm vào reward một số hạng $F=\gamma\Phi(s')-\Phi(s)$ (shaping theo
thế năng) thì policy tối ưu không đổi. Đây là cơ sở cho thành phần tiến độ trong reward của phần Improve.
\item \textbf{LIBERO} 130 task trong bốn bộ: Spatial, Object, Goal (mỗi bộ 10 task) và LIBERO-100
(LIBERO-90 + LIBERO-10/Long); mỗi task 50 demo do người điều khiển. Báo cáo dùng LIBERO-10 (Long).
\end{itemize}

\section*{Tài liệu tham khảo}
{\scriptsize\raggedright
\begin{itemize}
\item SimpleVLA-RL (2025). arXiv:2509.09674.
\item $\pi_{\text{RL}}$: Online RL fine-tuning for flow-based VLA. arXiv:2510.25889.
\item $\pi^*_{0.6}$ / RECAP. arXiv:2511.14759.
\item RIPT-VLA (2025). arXiv:2505.17016.
\item VLA-RL (2025). arXiv:2505.18719.
\item Diffusion Policy. arXiv:2303.04137.
\item $\pi_0$. arXiv:2410.24164.
\item SmolVLA. arXiv:2506.01844.
\item Flow Matching for Generative Modeling. arXiv:2210.02747.
\item Flow Matching Policy Gradients (FPO). arXiv:2507.21053.
\item Reinforcement Fine-Tuning of Flow-Matching Policies for VLA Models ($\pi_0$-FPO). arXiv:2510.09976.
\item Steering diffusion policy with latent space RL (DSRL). arXiv:2506.15799.
\item Policy invariance under reward transformations: theory and application to reward shaping. ICML.
\item LIBERO. NeurIPS Datasets and Benchmarks. arXiv:2306.03310.
\end{itemize}}

\end{document}
